\textbf{Задача 1.} Докажите, что язык $\mathrm{PAL}=\{x\in\{0,1\}^* | x^R=x\}$ лежит в классе $\dtime{n}$, но не распознается за время $O(n)$ на одноленточной машине Тьюринга.

\textbf{Интуиция:} На самом деле тут есть интуиция из теории информации, которую обычно явно не проговаривают, чтобы не уходить в изложении слишком далеко от основной темы.

Суть в том, что двум агентам для установления того факта, что им даны одинаковые строки длины n, нужно обменяться объёмом информации в хотя бы n битов (можно доказать по принципу Дирихле).

Разрезая ленту на две части, получаем двух агентов («Алисой» будет совокупность вычислений в одной части ленты, «Бобом» — совокупность вычислений в другой), которые посылают друг другу инфу в виде состояния машины Тьюринга.
Но состояний конечно, значит, машина перевозит только O(1) битов за раз, и потому ей нужно пересечь разрез хотя бы O(n) раз.

Если теперь сделать зону линейного размера, в каждой точке которой можно сделать разрез, для которого эти рассуждения работают, то получится линейное число участков, каждый из которых машина должна пересечь линейное число раз.

\begin{proof} \par
1. Легко заметить, что $\mathrm{PAL}$ лежит в $\dtime{n}$. Действительно, на двухленточной машине Тьюринга можно скопировать входную строку на вторую ленту, переместить на ней головку в конец строки, и, двигая головку на первой ленте вправо, а на второй влево, сравнить символы. Если всё совпало, то возвращаем 1, иначе 0.

2. Теперь покажем, что на одноленточной не может работать за время $O(n)$.

Рассмотрим слова длины $4n$ вида:
\[W_{x,y}=x\ 0^{2n}y^R, \quad |x|=|y|=n\]

Значит, $W_{x,y}$ палиндром $\iff$ $x=y$.

На каждой границе $i\in[n,3n]$ будем писать лог вида:
\[(state_1, direction_1), (state_2, direction_2), \ldots\]

То есть последовательность состояний и направлений (множество всех переходов) из $i-\text{ой}$ ячейки в $i+1-\text{ую}$.

\textbf{Утв.}
$\forall i\in[n,3n]\ \forall x\neq y$ логи на $i-\text{ой}$ границе слов $W_{x,x}$ и $W_{y,y}$ различны.

\begin{proof}
Иначе у машины будет одинаковое поведение на словах $W_{x,x}$ $W_{x,y}$, что не возможно, т.к для этих слов ответ разный. 
\end{proof}

Пусть алфавит для лога --- $Q$, $|Q| = q$.
Пусть $f(k) = (2n+1)\cdot(1+q+q^2+\cdots+q^k) = (2n+1)\cdot\frac{q^{2k+1}-1}{q-1}$ --- оценка сверху на максимальное количество различных слов (так как по утверждению один набор логов соответствует одному слову) длины $n$, у которых есть хотя бы один лог на позиции от $n$ до $3n$ длины не больше $k$. (на фикс. позиции $i$, таких логов максимум $(1+q+q^2+\cdots+q^k)$ (взятие $\le k$ элементов с возвращением)).

Возьмем такое $k$, что $f(k) < 2^n$, но $f(2k) \ge 2^n$.
Из первого неравенства получаем, что есть слово $x\in\{0,1\}^*\ |x|=n$ у которого все логи на позициях от $n$ до $3n$ длины больше $k$.

Из второго неравенства
\begin{align*}
f(2k)&=(2n+1)\cdot\frac{q^{2k+1}-1}{q-1}\ge 2n\cdot\frac{q^{2k+1}-1}{q-1}\ge 2^n\\
q^{2k+1} &\ge \frac{2^n}{2n}\cdot (q-1) \ge 2 ^{n/2}\cdot q 
\end{align*} (последнее неравенство верно, начиная с некоторого $n$)
\begin{align*}
	q^{2k}&\ge 2^{n/2}\\
	2k &\ge \frac{n/2}{\log_2 q}\\
	k &\ge \frac{n}{4\log_2 q}
\end{align*}

Значит, суммарная длина логов у слова $x$ будет больше или равна чем:
\[(2n + 1) \cdot k \ge 2n \cdot k \ge 2n \cdot \frac{n}{4\log_2 q} = \frac{n^2}{2\log_2 q} = O(n^2)\]

Поэтому время работы будет хотя бы квадратичным.
Что и требовалось доказать. 
\end{proof}
